\documentclass[10pt,a4paper]{amsart}
\usepackage[margin=19mm]{geometry}
\usepackage{amsmath,amssymb,mathtools}
\usepackage[scr=rsfs,cal=euler]{mathalfa}
\usepackage{xfrac}
\usepackage[unicode,hidelinks,bookmarks=false]{hyperref}
\newcommand{\ie}{i.e.\ }
\newcommand{\cf}{cf.\ }
\newcommand{\resp}{respectively\ }
\newcommand{\Subsection}{\subsection}
\providecommand{\nobreakdash}{\nobreak\hbox{-}}
\setlength{\emergencystretch}{2em}
\multlinegap0pt
\usepackage{bm}
\usepackage{bbm}
\usepackage{dsfont}
\usepackage{xspace}
\usepackage{cancel}
\usepackage{mathtools}
\usepackage{amsthm}
\makeatletter
\@ifundefined{definition}{\newtheorem{definition}{Definition}}{}
\@ifundefined{claimproof}{\newtheorem*{claimproof}{Claim}}{}
\@ifundefined{proposition}{\newtheorem{proposition}[definition]{Proposition}}{}
\makeatother
\renewcommand{\P}{\mathbb{P}}
\title[On the density of strongly minimal algebraic vector fields]{On the density of strongly minimal algebraic vector fields}
\author{R\'emi Jaoui}
\date{Source: arXiv:2301.06362v2 (CC BY 4.0)}
\begin{document}
\maketitle

\noindent\emph{Source and licence.} On the density of strongly minimal algebraic vector fields. Version arXiv:2301.06362v2 (CC BY 4.0), distributed under the Creative Commons Attribution 4.0 International licence, as recorded in the arXiv licence metadata for this exact version and shown on the abstract page https://arxiv.org/abs/2301.06362v2. Full rights notice: licenses/arxiv-2301.06362-CC-BY-4.0.txt.\par\smallskip

\noindent\textit{Editorial scope.} Counted unit: the Proposition whose source label is \texttt{presentation- twisted-vector field-hyperplanesection} in the article (source0.tex lines 532--536, source label \texttt{presentation- twisted-vector field-hyperplanesection}), delivered together with the article's own complete original proof of it (source0.tex lines 538--571, one \texttt{proof} environment of 34 lines). No mathematics is written, repaired or re-derived: statement and proof are the article's own text line for line. Required context retained uncounted: none. Every definition, display and label of those lines is retained unchanged. Omitted from this package and not redistributed: 1--531, 537--537, 572--2433. Nothing inside a retained excerpt is omitted or altered: every retained body is delivered exactly as the article's lines are written, so no omission receipt of any kind accompanies this package. Citations: the retained excerpts carry no citation command at all, so no bibliography entry of the article is transcribed and no reference is redistributed. Reference adaptation: none was needed and none was made; every reference inside the delivered unit resolves inside it, so no unresolved reference or citation is produced. Printed slips kept literal: no printed word, symbol, hypothesis or number of the counted statement or of its proof was altered. Layout and class: the article's own class is \texttt{amsart} with packages including amsmath, amssymb, amsfonts, graphics, natbib, babel, tikz, inputenc, xy, fontenc, pifont, hyperref, geometry, mathrsfs; the wrapper is the lane's pinned \texttt{amsart} editorial wrapper at 10pt on A4, which restates the theorem-like environments the retained text needs and adds the article's own macro definitions that the retained text needs (\texttt{\textbackslash P}), with extra wrapper packages (bm, bbm, dsfont, xspace, cancel, mathtools). The delivered numbering is the unit's own. Assets: no retained span contains a figure environment, a graphics-inclusion command, a PDF-page inclusion, a table or a TikZ picture; no image file is copied into this package, the package declares no asset dependency and no image file is redistributed with it. Licence: version arXiv:2301.06362v2 of this article is distributed under the Creative Commons Attribution 4.0 International licence, as recorded in the arXiv licence metadata for this exact version and shown on the abstract page https://arxiv.org/abs/2301.06362v2. A full rights notice is delivered at licenses/arxiv-2301.06362-CC-BY-4.0.txt. Characters: every retained excerpt is pure ASCII, so the delivered engine, XeLaTeX with its default Latin Modern text font, reports no missing character.
\begin{proposition} \label{presentation- twisted-vector field-hyperplanesection}
Let $X$ be a smooth  subvariety of the projective space $\P^N$, let $H_X = H \cap X$ be a smooth hyperplane section of $X$ and denote by $X_0$ be the complement of $H_X$ in $X$. The image of the restriction map 
$$-_{\mid X_0}: \Xi(X,\mathcal O_X(dH_X)) \rightarrow  \Xi(X_0) $$
is the $k$-vector space $\Xi(X_0,d)$ of regular vector fields on $X_0$ with a pole of order at most $d$ along the hypersurface at infinity $H_X$.  
\end{proposition} 

\begin{proof}
Denote by $\mathcal L = \mathcal O_X(H_X)$ so that $\mathcal L^{\otimes d} = \mathcal O_X(dH_X)$, fix a regular section $s$ of $\mathcal L$ on $X$ such that $\mathrm{div}(s) = H_X$ and define the restriction map
$$-_{\mid U_0}: \Xi(X,\mathcal L^{\otimes d}) \rightarrow \Xi(X_0)$$
relatively to the section $s^{\otimes d}$ of $\mathcal L^{\otimes d}$. Consider $\mathcal U = \lbrace U_0,\ldots, U_l \rbrace$ an affine open cover of $X$ such that each $U_i$ is equipped with a nowhere vanishing section $s_i$ of $\mathcal L_{\mid U_i}$. We may assume that $U_0 = X_0$ and that $s_0$ is the restriction of $s$ to $X_0$. By definition of the divisor of a section of line bundle, for every $i \leq l$, we can write 
$$ s_{\mid U_i} = h_i \cdot s_i$$ 
where $h_i \in \mathcal O_X(U_i)$ vanishes at order one on $H_X$ and hence is a local equation for $H_X$ on $U_i$ (so in particular $h_0 = 1$). It follows that the transition functions of $\mathcal L$ are given by 
$$\phi_{ij} = \frac {h_{i \mid U_i \cap U_j}} {h_{j \mid U_i \cap U_j}} \in \mathcal O^\ast _X(U_i \cap U_j).$$
They are indeed invertible since both $h_i$ and $h_j$ vanish only on $H$ and with the same order. Hence, an $\mathcal L^{\otimes d}$-twisted vector field can be identified with a collection 
$$ v:  \begin{cases} \lbrace (U_i,v_i) \mid i = 0,\ldots, l \rbrace \text{ of regular vector fields on each }U_i \\ 
\text{ such that } v_i = \phi_{ij}^d \cdot v_j \text{ on } U_i \cap U_j. \end{cases}
$$
We first show that the image of the restriction map is contained in $\Xi(X_0,d)$. Consider $v \in \Xi(X, \mathcal L^{\otimes d})$ as described above and consider $i$ such that $U_i \cap H \neq \emptyset$. On the open set $U_i \cap U_0$ we have 
$$h_i^d \cdot v_0 = \frac {h_i^d} {h_0^d} \cdot v_0 = \phi_{i0}^d \cdot v_0 = v_i$$
which is a regular vector field on $U_i$. It follows that $v_0$ which is the restriction of $v$ to $U$ has a pole of order at most $d$ along $H_X$ as required. To prove that the restriction map is surjective on $\Xi(X_0,d)$, we need the following claim: 

\begin{claimproof} 
Let $U$ be an affine open set of $X$ and let $H_U$ be an hypersurface of $U$. Assume that $v$ is a regular vector field on $U^\ast = U \setminus H_U$ with a pole of order at most zero along $H_U$. Then $v$ extends to a regular vector field $\overline{v}$ on $U$. 
\end{claimproof}

\begin{proof}[Proof of the claim] 
Denote by $\Theta_X(U)$ the free $\mathcal O_X(U)$-module of vector fields on $U$ and by $h$ an equation for $H_U$. Since $\Theta_X$ is a coherent sheaf, we have that $\Theta(U^\ast) = \Theta(U)_h$
is the localization of $\Theta(U)$ by the multiplicative set generated by $h$. It follows that 
$$h^n \cdot v  = v_1 \in \Theta_X(U)$$ 
extends to a regular vector field on $U$. Similarly, since $v$ has a pole of order at most zero along $H$, there is an irreducible $g \in \mathcal O_X(U)$  not associated with $h$ such that 
$$ g^m \cdot v = v_2 \in \Theta_X(U)$$
extends to a regular vector field on $U$ and we conclude that
$g^m\cdot v_1 = h^n \cdot v_2$.
Since  $X$ is smooth and $U$ is affine, $\mathcal O_X(U)$ is a unique factorization domain and we obtain that $g^m$ divides $v_2$ and that $h^n$ divides $v_1$ in $\Theta(U)$. This precisely means that $v$ extends to a vector field $\overline{v}$ on $U$. 
\end{proof} 
To finish the proof of surjectivity, consider $v_0 \in \Xi(X_0,d)$. 
For every $i = 1,\ldots, l$, the vector field $v_i = h_i^d \cdot v_0$ is a regular vector field on $U_i^\ast = U_i \setminus H_X$ with a pole of order at most zero along $H_X$. By the claim above, $v_i$ extends to a regular vector field $\overline{v_i}$ on $U_i$. It is now easy to see that the collection $(U_i, \overline{v_i})$ that we have constructed defines an $\mathcal L^{\otimes d}$-twisted vector field with restriction $v_0$: indeed, on $U_i \cap U_j \cap U_0$, we have
$$ \phi_{ij} \cdot \overline{v_j} = \frac{h_i^d}{h_j^d} \cdot \overline{v_j} =  h_i^d \cdot v_0 = \overline{v_i}$$  
and therefore the same equality holds in $U_i \cap U_j$. This finishes the proof of the proposition.
\end{proof}

\end{document}
